\documentclass[10pt]{article}

\usepackage[T1]{fontenc}
\usepackage[utf8]{inputenc}
\usepackage{lmodern}
\usepackage[a4paper,margin=0.78in,headheight=14pt]{geometry}
\usepackage{microtype}
\usepackage{amsmath,amssymb}
\usepackage{booktabs,longtable,tabularx,array}
\usepackage[table]{xcolor}
\usepackage{enumitem}
\usepackage{fancyhdr}
\usepackage{hyperref}
\usepackage{url}

\definecolor{navy}{HTML}{17365D}
\definecolor{blue}{HTML}{2F75B5}
\definecolor{pale}{HTML}{EAF2F8}
\definecolor{softgray}{HTML}{F3F4F6}
\definecolor{good}{HTML}{166534}
\definecolor{warn}{HTML}{92400E}

\hypersetup{
  colorlinks=true,
  linkcolor=navy,
  citecolor=blue,
  urlcolor=blue,
  pdftitle={VLM-as-a-Judge Baselines: Design, Motivation, and Results},
  pdfauthor={Ridham Patel}
}

\pagestyle{fancy}
\fancyhf{}
\lhead{VLM-as-a-Judge Baselines}
\rhead{Ridham Patel}
\cfoot{\thepage}
\renewcommand{\headrulewidth}{0.3pt}

\setlength{\parindent}{0pt}
\setlength{\parskip}{5pt}
\setlist[itemize]{leftmargin=1.5em,itemsep=2pt,topsep=3pt}
\setlist[enumerate]{leftmargin=1.7em,itemsep=2pt,topsep=3pt}
\DeclareRobustCommand{\method}[1]{\nolinkurl{#1}}
\newcommand{\target}{85\%}
\newcommand{\deltafull}{0.10}
\newcommand{\yes}{\textcolor{good}{yes}}
\newcommand{\no}{\textcolor{warn}{no}}
\newcolumntype{L}[1]{>{\raggedright\arraybackslash}p{#1}}

\title{\vspace{-1.2cm}\textcolor{navy}{\rule{\linewidth}{1.2pt}}\\[0.45cm]
\Huge\bfseries VLM-as-a-Judge Baselines\\[0.15cm]
\Large Design, Motivation, What Each Comparison Proves, and Results\\[0.45cm]
\textcolor{navy}{\rule{\linewidth}{1.2pt}}}
\author{Ridham Patel\\\href{mailto:ridham.patel@iitgn.ac.in}{ridham.patel@iitgn.ac.in}}
\date{October 2026}

\begin{document}
\maketitle

\begin{abstract}
This report documents every baseline and comparison policy implemented in the
multimodal selective-evaluation experiments. The task is pairwise VLM judging:
given an image, an instruction, and two candidate responses, a judge predicts the
response preferred by human annotators. The benchmark compares four model families,
four confidence-selection strategies, four prompt/order reliability setups, a
consistency-escalation policy, the original stage-wise certified cascade, and two
end-to-end certified routers. For each comparison, the report states the decision
rule, why the baseline is necessary, and the hypothesis it tests. It also records the
corrected 1,000-split four-judge results and identifies implemented policies whose
post-fix results have not yet been generated.
\end{abstract}

\tableofcontents
\newpage

\section{Experimental question and protocol}

The central question is not only whether a VLM can agree with humans. It is whether a
system can answer a useful fraction of comparisons while guaranteeing a user-selected
level of human agreement. The experiment sets the maximum selective error to
$\alpha=0.15$, so the accuracy target is $1-\alpha=\target$. The certificate uses
$\delta=\deltafull$, corresponding to a target confidence level of $1-\delta=90\%$.

The fast benchmark contains 1,000 approximately source-stratified, image-grounded
preference pairs. Each scored example contains three simulated-annotator prompts, one
few-shot demonstration per prompt, and both response orderings. The judge chain is:

\begin{center}
\fcolorbox{navy}{pale}{\parbox{0.92\linewidth}{\centering
Granite Vision 3.3 2B $\rightarrow$ Qwen2.5-VL 3B Instruct $\rightarrow$
Phi-3.5 Vision Instruct $\rightarrow$ InternVL3 8B}}
\end{center}

The chain deliberately spans different model families. This reduces the chance that a
result is merely a within-family scaling artifact and lets disagreement between model
families serve as a reliability signal. All policies operate on the same cached model
probabilities. Re-running policy evaluation requires no VLM inference.

\subsection{Metrics}

Let $A$ be the event that a policy answers and let $\widehat{Y}$ and $Y$ denote the
policy prediction and human label. The primary selective metric is
\[
  \operatorname{Accuracy}=\Pr(\widehat{Y}=Y\mid A),
\]
which is undefined for an empty selection and is therefore averaged only over nonempty
test splits. Coverage is $\Pr(A)$. Macro accuracy first computes accuracy within each
data source and then gives every source equal weight. Nonempty rate is the fraction of
repeated splits that answer at least one example. Target success is the fraction of all
splits that both answer and attain at least \target{} test accuracy. For heuristic and
point-estimate methods, ``success'' is empirical and is not a formal certificate.

\section{Why this baseline suite is necessary}

A single accuracy comparison cannot locate the source of improvement. The suite is
organized around five distinct questions:

\begin{longtable}{@{}p{0.20\linewidth}p{0.35\linewidth}p{0.38\linewidth}@{}}
\toprule
\textbf{Axis} & \textbf{Required comparison} & \textbf{What it identifies} \\
\midrule
\endhead
Judge strength & Direct use of every model & Whether gains come from a stronger judge rather than selection or routing. \\
Confidence selection & Heuristic, point estimate, and certified selectors & Whether confidence ranks correctness, and what finite-sample certification costs in coverage. \\
Cascading & Single-judge selectors versus sequential policies & Whether multiple model families add coverage beyond the strongest judge and what escalation costs. \\
Prompt reliability & Single prompt, order averaging, voting, and consistency filtering & Whether prompt variation and position bias explain confidence and accuracy. \\
Certificate scope & Stage-wise cascade versus global routers & Whether coverage is lost by protecting several stage claims instead of the deployed policy's final risk. \\
\bottomrule
\end{longtable}

The direct pairwise protocol follows multimodal judge and reward-model benchmarks that
measure agreement with human preferences \cite{chen2024mllmjudge,yasunaga2025mmrb,li2024vlreward}.
The selective and cascaded comparisons follow \emph{Trust or Escalate}
\cite{jung2024trust}. Prompt/order ablations operationalize reliability mechanisms used
in multimodal reward benchmarking and repeated judging. They are labeled as operational
analogues when the cached deterministic runs do not exactly reproduce a paper's original
stochastic protocol.

\section{Model-strength baselines}

\subsection[Direct judge]{Direct judge: \method{direct:<judge>}}

Each of the four direct baselines chooses $\arg\max_y p_j(y\mid x)$ and never abstains.
There is one instance for each model family; Appendix~A lists every exact identifier.

\textbf{Motivation.} Pairwise accuracy at 100\% coverage is the standard
VLM-as-a-judge comparison. It establishes the raw ability of each family and provides
the denominator against which abstention gains are judged.

\textbf{What it proves.} If a selective policy improves conditional accuracy, the direct
rows show how much is attributable to abstaining. Comparing the four direct rows also
tests whether the final judge is genuinely stronger. Direct baselines cannot demonstrate
risk control because they make no coverage-dependent decision and use no certificate.

\section{Confidence and selective-risk baselines}

\subsection[Strongest-judge heuristic]{Strongest-judge heuristic: \method{heuristic:internvl3-8b}}

The strongest judge answers when its cached confidence is at least $1-\alpha=0.85$.

\textbf{Motivation.} This is the simplest deployment rule a practitioner might use if
model probabilities were calibrated. It has no learned threshold and no statistical
bound.

\textbf{What it proves.} It tests whether confidence has useful ranking power and gives
an empirical upper reference for coverage. A high success rate does not turn it into a
certificate; it may fail under another calibration draw or distribution.

\subsection[Cascaded heuristic]{Cascaded heuristic: \method{cascaded_heuristic}}

The same fixed threshold, $0.85$, is applied sequentially. The first judge above the
threshold answers; otherwise the item is escalated.

\textbf{Motivation.} This is the non-certified counterpart of a cascade. It tests whether
sequential access to several families adds coverage beyond the strongest judge.

\textbf{What it proves.} Its difference from the strongest-judge heuristic isolates the
empirical value of escalation. Its difference from the certified cascade quantifies the
coverage price of finite-sample guarantees.

\subsection[Point-estimate calibration]{Point-estimate calibration: \method{point_estimate:<judge>}}

For each judge, the method chooses the least restrictive threshold whose empirical
calibration error is at most $\alpha$, requiring at least 30 selected calibration
examples. It does not add an upper confidence bound. Four instances are evaluated, one
for each model family.

\textbf{Motivation.} This is a calibrated but non-certified comparator from the selective
evaluation design. It separates threshold learning from statistical uncertainty.

\textbf{What it proves.} The gap between point-estimate and certified selection measures
the coverage cost of accounting for finite calibration size. Failure on held-out splits
shows why empirical calibration error alone is insufficient for a guarantee.

\subsection[Certified single judge]{Certified single judge: \method{selective:<judge>}}

Each judge independently uses fixed-sequence testing and the full $\delta=0.1$ to find a
threshold whose one-sided exact binomial upper risk bound is at most $\alpha$. There is one
instance for each model family; Appendix~A lists every exact identifier.

\textbf{Motivation.} A single certified judge is the essential control for the cascade.
It receives the full failure budget and does not benefit from other models.

\textbf{What it proves.} It separates gains from cascading from losses caused by splitting
$\delta$ across stages. These policies are implemented in the current code, but they were
added after the captured result file was generated. Their rows are therefore marked
``N/R'' rather than filled with inferred values. Re-running CPU evaluation will produce
them from the existing caches.

\subsection[Original certified cascade]{Original certified cascade: \method{cascaded_selective}}

At stage $j$, fixed-sequence testing certifies a confidence threshold on the residual
examples rejected by earlier stages. With four judges, the implementation assigns
$\delta_j=\delta/4=0.025$ to each stage. If every stage-specific risk is at most $\alpha$,
the final risk is at most $\alpha$ with probability at least $1-\delta$ by the union bound.

\textbf{Motivation.} This is the original proposed method and the principal comparator
for the new global certificate \cite{jung2024trust}.

\textbf{What it proves.} It tests whether a formally protected cascade can meet the human
agreement target. Its coverage reveals the conservatism of simultaneously certifying
multiple stages on shrinking residual sets.

\section{Prompt and response-order reliability baselines}

All four setup policies are instantiated separately for Granite, Qwen, Phi, and InternVL.
They are post-hoc transformations of cached per-prompt, per-order probabilities, so their
comparison requires no new model calls.

\subsection[Vanilla single prompt]{Vanilla single prompt: \method{setup:vanilla_single:<judge>}}

Use the first simulated-annotator prompt and the original response order.

\textbf{Motivation.} This reproduces the common fixed-prompt judge setup used in basic
pairwise evaluation \cite{chen2024mllmjudge,yasunaga2025mmrb}.

\textbf{What it proves.} It establishes how much the full prompt/order ensemble gains over
one ordinary judge call. It also exposes position sensitivity because it never reverses
the candidates.

\subsection[Response-order averaging]{Response-order averaging: \method{setup:position_swap:<judge>}}

Evaluate the first prompt in both response orders, map both outputs back to the original
candidate identities, average their probabilities, and always answer.

\textbf{Motivation.} Pairwise judges can prefer the first or second slot. Forward/reverse
evaluation is used in multimodal reward benchmarking to measure and reduce this bias
\cite{hu2025mmrb2}.

\textbf{What it proves.} Improvement over the vanilla row measures the benefit of explicit
position control. Because it retains full coverage, it does not confound the comparison
with abstention.

\subsection[Prompt/order vote]{Prompt/order vote: \method{setup:prompt_order_vote:<judge>}}

Take a majority vote across every cached simulated-annotator prompt and both response
orders.

\textbf{Motivation.} Repeated judgments are a standard way to reduce prompt variance;
VLRewardBench studies inference-time scaling for generative reward models
\cite{li2024vlreward}.

\textbf{What it proves.} It tests whether diversity across prompts and response orders
provides a robust ensemble signal. This implementation is a deterministic analogue over
cached runs, not a reproduction of a stochastic five-call protocol.

\subsection[Swap consistency]{Swap consistency: \method{setup:swap_consistency:<judge>}}

Answer only when the first prompt selects the same original response in both response
orders; otherwise abstain. For consistent pairs, average the two probability vectors.

\textbf{Motivation.} Self-consistency can serve as a label-free uncertainty signal and
directly targets position-instability failures.

\textbf{What it proves.} Its accuracy/coverage curve tests whether order agreement
identifies reliable examples. Comparing it with position averaging isolates the effect of
filtering inconsistent items rather than merely averaging them.

\subsection[Consistency escalation]{Consistency escalation: \method{setup:consistency_escalation}}

At each non-final stage, accept only swap-consistent judgments and route inconsistent
examples onward. The final judge resolves every remaining pair using its full cached
ensemble.

\textbf{Motivation.} Selective test-time evaluation allocates extra computation to
self-inconsistent cases \cite{jwa2026lwe}. This policy tests that reusable routing idea
without implementing an evolving meta-prompt.

\textbf{What it proves.} It tests whether consistency alone is sufficient for economical
escalation. It is an operational analogue of the routing heuristic, not a full
implementation of Learning While Evaluating.

\section{End-to-end certified routing}

\subsection{Why change the scope of certification?}

The stage-wise guarantee is valid, but it controls a stronger set of events than the
deployed system requires. It separately requires
\[
R_j=\Pr(\widehat{Y}_j\ne Y\mid\text{reached and accepted at stage }j)\leq\alpha
\]
for every stage. The union bound protects all $M$ claims with
$\sum_j\delta_j=\delta$. This is conservative because each test uses a smaller $\delta_j$,
later stages receive fewer calibration examples, and every stage must meet $\alpha$ even
though final risk is a coverage-weighted mixture of stage errors.

The global experiment freezes the complete policy $\pi_\tau$ and directly certifies
\[
R(\pi_\tau)=\Pr(\widehat{Y}_{\pi_\tau}\ne Y\mid\pi_\tau\text{ answers})\leq\alpha.
\]
One policy-level claim can use the full $\delta$. To prevent label leakage, every
repetition uses disjoint group-safe router-training, certification, and test partitions.
Threshold candidates are fixed from router-training scores before certification labels are
examined.

\subsection[Global raw confidence]{Global raw confidence: \method{global_raw_confidence}}

Route using each stage's cached maximum class probability, freeze a candidate threshold
grid from the router-training split, and certify the final routing policy once.

\textbf{Motivation.} This is the minimal end-to-end modification. It changes certificate
scope while keeping the original confidence signal unchanged.

\textbf{What it proves.} Its comparison with \method{cascaded_selective} isolates the
benefit of a policy-level certificate. Its comparison with the learned router isolates
whether richer reliability features help beyond raw confidence.

\subsection[Global learned router]{Global learned router: \method{global_learned_router}}

Fit one regularized logistic correctness model per stage on the independent router split.
Features are confidence, margin, entropy, simulated-annotator vote support, probability
dispersion, response-order consistency and probability shift, agreement/unanimity among
models visited so far, and previous-confidence summaries. The learned score is then frozen
and certified exactly like the raw score.

\textbf{Motivation.} VLM reliability may be expressed by disagreement and response-order
instability even when maximum probability is poorly calibrated.

\textbf{What it proves.} It tests whether cached multimodal reliability signals can improve
routing without new inference. A loss against raw confidence is informative evidence that
the 11-feature model has too much variance for the approximately 200-example router split.

\section{Complete four-judge results}

Table~\ref{tab:baseline-results} reproduces the corrected post-fix baseline output. Accuracy
excludes empty selections. The four certified single-judge policies are present in the
current experiment registry but absent from the captured run, so no numbers are imputed.

\footnotesize
\begin{longtable}{@{}L{0.50\linewidth}rrrrr@{}}
\caption{Four-judge baseline results over 1,000 repeated group-safe splits. Values are
fractions. N/R means implemented but not reported in the captured run.}\label{tab:baseline-results}\\
\toprule
\textbf{Policy} & \textbf{Acc.} & \textbf{Macro} & \textbf{Cov.} & \textbf{Nonempty} & \textbf{Success} \\
\midrule
\endfirsthead
\toprule
\textbf{Policy} & \textbf{Acc.} & \textbf{Macro} & \textbf{Cov.} & \textbf{Nonempty} & \textbf{Success} \\
\midrule
\endhead
\method{direct:granite-vision-3.3-2b} & .5909 & .5971 & 1.0000 & 1.0000 & .0000 \\
\method{direct:qwen2.5-vl-3b-instruct} & .5762 & .5988 & 1.0000 & 1.0000 & .0000 \\
\method{direct:phi-3.5-vision-instruct} & .5824 & .6115 & 1.0000 & 1.0000 & .0000 \\
\method{direct:internvl3-8b} & .7051 & .7325 & 1.0000 & 1.0000 & .0000 \\
\rowcolor{pale}\method{heuristic:internvl3-8b} & .8872 & .8573 & .3380 & 1.0000 & .9990 \\
\rowcolor{pale}\method{cascaded_heuristic} & .8695 & .8477 & .3611 & 1.0000 & .9660 \\
\method{point_estimate:granite-vision-3.3-2b} & .6816 & .6931 & .0099 & .0580 & .0000 \\
\method{point_estimate:qwen2.5-vl-3b-instruct} & .8347 & .7792 & .1320 & .8650 & .3250 \\
\method{point_estimate:phi-3.5-vision-instruct} & .7264 & .7102 & .0240 & .1770 & .0000 \\
\method{point_estimate:internvl3-8b} & .8477 & .8368 & .4274 & 1.0000 & .4570 \\
\method{selective:granite-vision-3.3-2b} & N/R & N/R & N/R & N/R & N/R \\
\method{selective:qwen2.5-vl-3b-instruct} & N/R & N/R & N/R & N/R & N/R \\
\method{selective:phi-3.5-vision-instruct} & N/R & N/R & N/R & N/R & N/R \\
\method{selective:internvl3-8b} & N/R & N/R & N/R & N/R & N/R \\
\rowcolor{pale}\method{cascaded_selective} & .9362 & .8793 & .1306 & .6680 & .6370 \\
\method{setup:vanilla_single:granite-vision-3.3-2b} & .4998 & .5016 & 1.0000 & 1.0000 & .0000 \\
\method{setup:position_swap:granite-vision-3.3-2b} & .6125 & .6394 & 1.0000 & 1.0000 & .0000 \\
\method{setup:prompt_order_vote:granite-vision-3.3-2b} & .5784 & .5953 & 1.0000 & 1.0000 & .0000 \\
\method{setup:swap_consistency:granite-vision-3.3-2b} & .7640 & .8096 & .0681 & 1.0000 & .0080 \\
\method{setup:vanilla_single:qwen2.5-vl-3b-instruct} & .5830 & .6045 & 1.0000 & 1.0000 & .0000 \\
\method{setup:position_swap:qwen2.5-vl-3b-instruct} & .5812 & .6060 & 1.0000 & 1.0000 & .0000 \\
\method{setup:prompt_order_vote:qwen2.5-vl-3b-instruct} & .5883 & .6072 & 1.0000 & 1.0000 & .0000 \\
\method{setup:swap_consistency:qwen2.5-vl-3b-instruct} & .6606 & .6225 & .4325 & 1.0000 & .0000 \\
\method{setup:vanilla_single:phi-3.5-vision-instruct} & .5641 & .5899 & 1.0000 & 1.0000 & .0000 \\
\method{setup:position_swap:phi-3.5-vision-instruct} & .5709 & .6025 & 1.0000 & 1.0000 & .0000 \\
\method{setup:prompt_order_vote:phi-3.5-vision-instruct} & .5772 & .6012 & 1.0000 & 1.0000 & .0000 \\
\method{setup:swap_consistency:phi-3.5-vision-instruct} & .6379 & .6628 & .4940 & 1.0000 & .0000 \\
\method{setup:vanilla_single:internvl3-8b} & .6631 & .6949 & 1.0000 & 1.0000 & .0000 \\
\method{setup:position_swap:internvl3-8b} & .7030 & .7289 & 1.0000 & 1.0000 & .0000 \\
\method{setup:prompt_order_vote:internvl3-8b} & .6960 & .7230 & 1.0000 & 1.0000 & .0000 \\
\method{setup:swap_consistency:internvl3-8b} & .7729 & .7830 & .6522 & 1.0000 & .0000 \\
\method{setup:consistency_escalation} & .6150 & .6430 & 1.0000 & 1.0000 & .0000 \\
\bottomrule
\end{longtable}
\normalsize

The two global policies use a three-way split (approximately 200 router-training, 300
certification, and 500 test examples), whereas the baseline table uses approximately 300
calibration and 700 test examples. The global policies are paired with each other on the
same splits; comparisons with Table~\ref{tab:baseline-results} are repeated-split aggregate
comparisons rather than paired test-set differences.

\begin{table}[ht]
\centering
\caption{End-to-end certified routing over 1,000 repeated splits.}
\begin{tabular}{@{}lrrrrrrr@{}}
\toprule
\textbf{Method} & \textbf{Acc.} & \textbf{Macro} & \textbf{Cov.} & \textbf{Calls} & \textbf{Nonempty} & \textbf{Certified} & \textbf{Success} \\
\midrule
\method{global_raw_confidence} & .9108 & .8658 & .2520 & 3.891 & .9600 & .9600 & .8650 \\
\method{global_learned_router} & .9042 & .8604 & .1543 & 3.909 & .6180 & .6180 & .5370 \\
\bottomrule
\end{tabular}
\end{table}

\section{What the results establish}

\begin{enumerate}
  \item \textbf{Judge strength matters, but is insufficient.} InternVL is the strongest
  direct model at 70.51\%, still well below the 85\% target. A stronger final model alone
  does not solve reliable judging.

  \item \textbf{Confidence is useful.} The InternVL heuristic reaches 88.72\% accuracy at
  33.80\% coverage. The point-estimate version reaches greater coverage (42.74\%) but only
  84.77\% mean accuracy and succeeds on 45.70\% of splits. This is direct evidence that
  uncertainty in threshold calibration matters.

  \item \textbf{Position control helps unevenly.} Granite improves from 49.98\% with one
  prompt/order to 61.25\% after order averaging. InternVL improves from 66.31\% to 70.30\%.
  Qwen and Phi change little. Position bias is therefore model dependent.

  \item \textbf{Consistency is selective but not sufficient.} Swap consistency raises
  conditional accuracy for every judge, but only Granite attains a small nonzero target
  success rate. The full-coverage consistency-escalation analogue reaches 61.50\%, showing
  that consistency alone is a weak routing signal in this chain.

  \item \textbf{Stage-wise certification is accurate but conservative.} The original
  cascade reaches 93.62\% accuracy, 8.62 points above target, at 13.06\% coverage and a
  66.80\% nonempty rate. This large unused accuracy margin motivates a less conservative
  policy-level certificate.

  \item \textbf{Global certification nearly doubles coverage.} Raw-confidence global
  certification raises coverage from 13.06\% to 25.20\% (1.93$\times$) while maintaining
  91.08\% accuracy. Among certified splits, target success is
  $0.865/0.960=90.1\%$, matching the intended 90\% confidence level empirically.

  \item \textbf{The rich router is a negative ablation.} The learned router has lower
  coverage and certifies fewer splits than raw confidence. With only about 200 router
  examples, the 11-feature correctness model likely has too much variance.

  \item \textbf{Compute efficiency remains unresolved.} The raw global policy calls 3.891
  of four judges on average. Its main benefit is coverage from certificate scope, not early
  exit. Future routing work should report judge calls alongside accuracy and coverage.
\end{enumerate}

\section{Completeness and scope boundary}

The experiment registry contains 34 distinct policy instances:

\begin{center}
\begin{tabular}{@{}lr@{}}
\toprule
\textbf{Policy family} & \textbf{Instances} \\
\midrule
Direct judges & 4 \\
Strongest-judge heuristic & 1 \\
Cascaded heuristic & 1 \\
Point-estimate single judges & 4 \\
Certified single judges & 4 \\
Original certified cascade & 1 \\
Vanilla single-prompt setups & 4 \\
Position-swap setups & 4 \\
Prompt/order-vote setups & 4 \\
Swap-consistency setups & 4 \\
Consistency escalation & 1 \\
Global raw-confidence router & 1 \\
Global learned router & 1 \\
\midrule
\textbf{Total} & \textbf{34} \\
\bottomrule
\end{tabular}
\end{center}

Prometheus-Vision \cite{lee2024prometheusvision} and LLaVA-Critic
\cite{xiong2024llavacritic} are relevant specialist evaluator models, but they are not rows
in the current cached experiment. Prometheus-Vision is trained for rubric-conditioned
pointwise scoring, whose primary metric is correlation; adapting it to pairwise preference
would change the protocol. LLaVA-Critic supports pairwise ranking, but a faithful comparison
requires its official generative prompt/parser and new cached inference. Neither is silently
presented as an implemented baseline. They should be added as separate model baselines only
after their native protocols are validated on the same examples.

Multimodal RewardBench, MMRB2, MLLM-as-a-Judge, and VLRewardBench motivate the pairwise
metrics and prompt/order controls; they are benchmark or setup references rather than
additional model rows. Learning While Evaluating motivates the consistency-escalation
heuristic; the current baseline does not claim its evolving meta-prompt mechanism.

\section{Reproduction and reporting checklist}

The complete cached-only evaluation command is:
\begin{verbatim}
python -u -m open_cascade evaluate configs/vlm_v2_fast_4judge.yaml
\end{verbatim}

For a paper table:
\begin{itemize}
  \item rerun the current code to populate the four \method{selective:<judge>} rows;
  \item retain 1,000 group-safe repetitions and report mean plus standard deviation;
  \item report accuracy, macro accuracy, coverage, macro coverage, nonempty rate, and target success;
  \item for global routing, also report certified rate, mean judge calls, early-exit rate, and model composition;
  \item state that raw/learned global comparisons are paired, while the original/global tables use different test sizes;
  \item label heuristic and point-estimate success as empirical rather than guaranteed;
  \item keep empty selections out of conditional-accuracy averages.
\end{itemize}

\appendix
\section{Exact implemented policy identifiers}

\footnotesize
\begin{longtable}{@{}L{0.23\linewidth}L{0.71\linewidth}@{}}
\toprule
\textbf{Family} & \textbf{Identifiers} \\
\midrule
\endhead
Direct & \method{direct:granite-vision-3.3-2b}; \method{direct:qwen2.5-vl-3b-instruct}; \method{direct:phi-3.5-vision-instruct}; \method{direct:internvl3-8b}. \\
Heuristic & \method{heuristic:internvl3-8b}; \method{cascaded_heuristic}. \\
Point estimate & \method{point_estimate:granite-vision-3.3-2b}; \method{point_estimate:qwen2.5-vl-3b-instruct}; \method{point_estimate:phi-3.5-vision-instruct}; \method{point_estimate:internvl3-8b}. \\
Certified single judge & \method{selective:granite-vision-3.3-2b}; \method{selective:qwen2.5-vl-3b-instruct}; \method{selective:phi-3.5-vision-instruct}; \method{selective:internvl3-8b}. \\
Certified cascade & \method{cascaded_selective}. \\
Vanilla setup & \method{setup:vanilla_single:granite-vision-3.3-2b}; \method{setup:vanilla_single:qwen2.5-vl-3b-instruct}; \method{setup:vanilla_single:phi-3.5-vision-instruct}; \method{setup:vanilla_single:internvl3-8b}. \\
Position swap & \method{setup:position_swap:granite-vision-3.3-2b}; \method{setup:position_swap:qwen2.5-vl-3b-instruct}; \method{setup:position_swap:phi-3.5-vision-instruct}; \method{setup:position_swap:internvl3-8b}. \\
Prompt/order vote & \method{setup:prompt_order_vote:granite-vision-3.3-2b}; \method{setup:prompt_order_vote:qwen2.5-vl-3b-instruct}; \method{setup:prompt_order_vote:phi-3.5-vision-instruct}; \method{setup:prompt_order_vote:internvl3-8b}. \\
Swap consistency & \method{setup:swap_consistency:granite-vision-3.3-2b}; \method{setup:swap_consistency:qwen2.5-vl-3b-instruct}; \method{setup:swap_consistency:phi-3.5-vision-instruct}; \method{setup:swap_consistency:internvl3-8b}. \\
Consistency escalation & \method{setup:consistency_escalation}. \\
Global routers & \method{global_raw_confidence}; \method{global_learned_router}. \\
\bottomrule
\end{longtable}
\normalsize

\begin{thebibliography}{99}

\bibitem{jung2024trust}
J. Jung, F. Brahman, and Y. Choi.
\newblock Trust or Escalate: LLM Judges with Provable Guarantees for Human Agreement.
\newblock \emph{arXiv:2407.18370}, 2024.

\bibitem{chen2024mllmjudge}
D. Chen et al.
\newblock MLLM-as-a-Judge: Assessing Multimodal LLM-as-a-Judge with Vision-Language Benchmark.
\newblock \emph{arXiv:2402.04788}, 2024.

\bibitem{yasunaga2025mmrb}
M. Yasunaga, L. Zettlemoyer, and M. Ghazvininejad.
\newblock Multimodal RewardBench: Holistic Evaluation of Reward Models for Vision Language Models.
\newblock \emph{arXiv:2502.14191}, 2025.

\bibitem{li2024vlreward}
L. Li et al.
\newblock VLRewardBench: A Challenging Benchmark for Vision-Language Generative Reward Models.
\newblock \emph{arXiv:2411.17451}, 2024.

\bibitem{hu2025mmrb2}
Y. Hu, R. Askari-Hemmat, M. Hall, E. Dinan, L. Zettlemoyer, and M. Ghazvininejad.
\newblock Multimodal RewardBench 2: Evaluating Omni Reward Models for Interleaved Text and Image.
\newblock \emph{arXiv:2512.16899}, 2025.

\bibitem{lee2024prometheusvision}
S. Lee, S. Kim, S. H. Park, G. Kim, and M. Seo.
\newblock Prometheus-Vision: Vision-Language Model as a Judge for Fine-Grained Evaluation.
\newblock \emph{arXiv:2401.06591}, 2024.

\bibitem{xiong2024llavacritic}
T. Xiong et al.
\newblock LLaVA-Critic: Learning to Evaluate Multimodal Models.
\newblock \emph{arXiv:2410.02712}, 2024.

\bibitem{jwa2026lwe}
S. Jwa, D. Ahn, R. Kim, D. Kang, and J. Choi.
\newblock Becoming Experienced Judges: Selective Test-Time Learning for Evaluators.
\newblock In \emph{Proceedings of EACL}, pages 697--721, 2026.

\end{thebibliography}

\end{document}
